
Reinforcement Learning from Human Feedback (RLHF) relies on human preference labels to train reward models, which are then used to fine-tune language models via policy gradient methods. A foundational assumption of this pipeline is that human annotators are stable, consistent oracles of response quality. Analysis of the Anthropic HH-RLHF dataset — 160,800 preference comparisons structured across three annotation phases — reveals a pattern that challenges this assumption: the coefficient associated with response verbosity in a logistic preference model reverses sign from early to late annotation strata. Early-stratum annotators exhibit a negative verbosity weight (β_L = −0.025), while late-stratum annotators exhibit a positive weight (β_L = +0.056). This shift is associated with annotation round and is directionally consistent with the stylistic profile of AI-generated text under annotation.

This observation raises the question of whether the human preference signal in multi-round RLHF annotation is stationary, and whether any directional shift is specifically aligned with AI-typical stylistic features rather than reflecting arbitrary variation or quality recalibration.

\textbf{The gap addressed.} Prior work on reward misspecification [Pan et al. 2022] and annotation variance [Coste et al. 2023] treats annotator preferences as stationary noise rather than as potentially directionally drifting signals. The temporal structure of multi-round annotation datasets such as HH-RLHF and WebGPT has not been exploited to analyze preference signal stationarity. No instrument exists for monitoring annotation-level Human→AI stylistic adaptation in RLHF pipelines.

\textbf{Approach.} This paper introduces the \textbf{Alignment Asymmetry Index (AAI)}, a composite instrument for measuring directional stylistic drift in RLHF preference annotations. The AAI is designed around three components: (1) directional stylistic coefficient drift (Δβ_L, Δβ_H, Δβ_S) across annotation strata after controlling for a quality covariate; (2) cosine projection of the annotation preference gradient onto a pre-defined AI-typicality embedding vector; and (3) behavioral divergence between reward models trained on early versus late annotation rounds. Components 1 and 2 are evaluated in this paper; component 3 is specified for future work.

\textbf{Contributions:}

\begin{enumerate}
\item \textbf{Verbosity-specific annotation drift measurement (H-M2).} The first coefficient-resolution measurement of verbosity preference reversal across annotation strata in a real RLHF dataset: early-stratum annotators are associated with penalizing verbosity, while late-stratum annotators are associated with rewarding it (Δβ_L = +0.080, non-overlapping 95\% bootstrap CIs; 2,000 stratified resamples; 53,600 rows per stratum). The hedging and structured reasoning coefficients show positive Δ but overlapping CIs, so the pre-registered gate of at least two non-overlapping features is not met (n_directional = 1 of 3).
\end{enumerate}

\begin{enumerate}
\item \textbf{AI-typicality geometric projection with discriminant validity (H-M1, AAI components 1–2).} A between-group tercile design on WebGPT yields a statistically significant tercile F-statistic (F = 33.226, p ≈ 8.3×10⁻⁹). A placebo permutation test on a random embedding direction yields an empirical p-value of 0.48, confirming that the between-group difference is specific to the AI-typicality direction. The pre-registered effect-size threshold (β_exposure ≥ 0.1 SD per 1,000 tokens) was not met; the executed β_exposure estimate is approximately 4.1×10⁻⁵ in the between-group regression, and the within-annotator panel design could not be executed due to the absence of worker identity metadata.
\end{enumerate}

\begin{enumerate}
\item \textbf{Minimum data requirements specification (H-E1, H-M1, H-M2).} Null results — interaction p = 1.0 (absent ambiguity labels), effect size below pre-registered threshold (absent within-annotator tracking), topic imbalance p = 4×10⁻²⁷⁵ — constitute a principled, empirically grounded specification of the data requirements for confirming individual annotator causal adaptation.
\end{enumerate}

